\section{A constructive causal acquired-geometry theorem}\label{sec:central}
We first define the finite-dimensional feasible object in the theorem. Fix live and terminal sets $S,J$ and $s_0\in S$. Put
\[
 a_s^\theta=\one_{\{s=s_0\}}x_0^\theta.
\]
A candidate consists of common action probabilities $p_{sa}$, common decoders $r_j\in H$, and, for each $\theta$, vectors $z_s^\theta,d_j^\theta\in C$ with
\begin{equation}\label{eq:budget-candidate}
 z_s^\theta-a_s^\theta\in C,\qquad
 \sum_s e(z_s^\theta)\le N,\qquad
 \sum_j e(d_j^\theta)=1.
\end{equation}
For every finite $F\subset\Theta$ and every collection $u_s^\theta,v_j^\theta\in E^*$, require
\begin{align}
 &\sum_{\theta\in F}\left[
   \sum_{s\in S}\ip{u_s^\theta}{z_s^\theta-a_s^\theta}
       +\sum_{j\in J}\ip{v_j^\theta}{d_j^\theta}\right]\notag\\
 &\quad\le\sum_{s\in S,\,a\in\calA}p_{sa}
    \int \max\left\{
      \max_{i\in S}\sum_{\theta\in F}\ip{u_i^\theta}{L_a^\theta(y)z_s^\theta},
      \max_{j\in J}\sum_{\theta\in F}\ip{v_j^\theta}{L_a^\theta(y)z_s^\theta}
              \right\}\lambda_a(dy).\label{eq:joint-support}
\end{align}
The summation over models precedes the maximum. The maximum chooses one next label, including the decision to stop, and cannot be chosen separately for different models. Define
\begin{equation}\label{eq:Phi}
 \Phi^\Theta(N,M)=
 \min\ \sup_{\theta\in\Theta}\sum_{j\in J}G_\theta(d_j^\theta,r_j),
\end{equation}
where the minimum is over the finitely many cardinalities $|S|,|J|\ge1$, $|S|+|J|\le M$, and over candidates satisfying \eqref{eq:budget-candidate}--\eqref{eq:joint-support}. The minimum and its operational meaning are part of the following result.

For the geometric part, let $K\subset\R^k$ be compact and $\nu$ a probability measure with support $K$. Write
\begin{equation}\label{eq:profile}
 \calQ_m^\nu(t)=\min_{c_1,\ldots,c_m\in\conv K}
 \ \min_{\substack{0\le h\le t\ (\nu\text{-a.e.})\\\int h\,d\nu=1}}
       \int_K\min_{1\le j\le m}\norm{x-c_j}^2h(x)\nu(dx),
 \qquad t\ge1.
\end{equation}
Both minima are attained, as proved below. The factor $t$ is a constraint on acquired probability mass, not a radius or a dimension.

\begin{definition}[Report-recoverable task with comparable acquisition]\label{def:recoverable}
A task is report-recoverable if, whenever action $a$ produces report $y$, its current conditional audit mean is a specified Borel function $x=f_a(y)\in K$, independent of the incoming history and of $\theta$, and its conditional audit variance is a common constant $B$. Suppose $\nu$ is the actual one-call law under a specified admissible reference preparation and action. There are constants $0<a_*\le1\le b_*<\infty$ such that, for every model and every positive-probability visible history, every legal next action has task-response law at most $b_*\nu$. One fixed available action $a_0$, common to all models, has conditional task-response law at least $a_*\nu$ at every such history. These are measure inequalities, not pointwise assertions on null histories.
\end{definition}
The lower inequality is needed only for the implementing action. A raw-kernel bound on a closed invariant base is one way, used in the realizations, to verify the required history quantifiers. The actual terminal law remains policy-dependent.

\begin{theorem}[Constructive budgeted causal acquired-geometry transfer]\label{thm:main}
The following statements hold.

\smallskip\noindent\textup{(i) General implementation and robust attainment.}
Under \eqref{eq:raw-instrument}--\eqref{eq:terminal-cost}, for every real $N\ge1$ and integer $M\ge2$,
\begin{equation}\label{eq:value-equality}
 \calR_\aut^\Theta(N,M)=\Phi^\Theta(N,M).
\end{equation}
A common Borel program attains this value. No uniform transience, positive stopping floor, mixing assumption, or finite bound on $\tau$ is imposed. In addition,
\begin{equation}\label{eq:finite-models}
 \calR_\aut^\Theta(N,M)
 =\sup_{\substack{F\subset\Theta\\0<|F|<\infty}}
         \calR_\aut^F(N,M).
\end{equation}
Thus every risk level strictly below the full-family optimum is already impossible on some finite model subfamily, with the same $N,M$ and one common program.

\smallskip\noindent\textup{(ii) Actual-law geometric bounds.}
For a report-recoverable task as in Definition~\ref{def:recoverable}, all fully adaptive checkpoint policies and all admissible autonomous programs satisfy
\begin{align}
 \calQ_m^\nu(b_*N)&\le\calR_\cp^\Theta(N,m)-B,
 \label{eq:profile-check-lower}\\
 \calQ_{M-1}^\nu(b_*N)&\le\calR_\aut^\Theta(N,M)-B.
 \label{eq:profile-aut-lower}
\end{align}
If $a_*N\ge1$, one common autonomous program gives
\begin{equation}\label{eq:profile-aut-upper}
 \calR_\aut^\Theta(N,M)-B
 \le\frac{b_*}{a_*}\calQ_{M-1}^\nu(a_*N).
\end{equation}
It has exactly one live state and at most $M-1$ terminal states, and its stopping time has a conditional geometric tail with mean at most $N$. The analogous checkpoint upper uses $m$ terminal states in place of $M-1$.

If every conditional response law equals $\nu$, independently of all previous reports and all actions, then for all $N\ge1$,
\begin{equation}\label{eq:exact-profile}
 \calR_\cp^\Theta(N,m)-B=\calQ_m^\nu(N),\qquad
 \calR_\aut^\Theta(N,M)-B=\calQ_{M-1}^\nu(N).
\end{equation}
The realizing rule in this case has $\E\tau=N$, including noninteger $N$.

\smallskip\noindent\textup{(iii) Matching resource law on acquired regular supports.}
Suppose in addition that, for some $s>0$, $r_0>0$ and $0<c_0\le C_0<\infty$,
\begin{equation}\label{eq:regularity}
 c_0r^s\le\nu(B(x,r))\le C_0r^s
 \quad(x\in K,\ 0<r\le r_0).
\end{equation}
Then, simultaneously for all $N\ge1$, $m\ge1$ and $M\ge2$,
\begin{align}
 \calR_\cp^\Theta(N,m)-B&\asymp (Nm)^{-2/s},\label{eq:regular-cp}\\
 \calR_\aut^\Theta(N,M)-B&\asymp \bigl(N(M-1)\bigr)^{-2/s}.\label{eq:regular-aut}
\end{align}
The constants depend only on $a_*,b_*,c_0,C_0,r_0,s$ and the diameter of $K$, not on $N,m,M$, the number of models, or the chosen adaptive policy.

\smallskip\noindent\textup{(iv) Finite construction at exactly the declared budgets.}
Under the effective raw-instrument certificate of Definition~\ref{def:effective}, for every rational $N\ge1$, integer $M\ge2$, and rational $\varepsilon>0$, a finite calculation returns a common finitely described observer and a certified risk interval of width at most $\varepsilon$. Its risk lies in the interval above $\Phi^\Theta(N,M)$, while the interval's lower endpoint bounds every admissible Borel observer. The returned observer still has mean acquisition at most $N$ and total state count at most $M$. Neither report recoverability nor posterior contraction is needed for this part. Stopping regularization and occupation-weighted error estimates, rather than a free deterministic clock or a budget relaxation, supply the construction.

\smallskip\noindent\textup{(v) Sharp-order transfer at unchanged resources.}
Consider two experiment families with the same declared interfaces and uniform preparation, joint one-call and audit errors $\delta_0,\eta,\alpha$ in the base-norm sense of Section~\ref{sec:sharp-transfer}. Put $\zeta=\delta_0+N\eta$ and $c_N=1-N^{-1}$. Their optimal values satisfy
\begin{equation}\label{eq:central-normalized}
 |\Phi_1^\Theta(N,M)-\Phi_2^\Theta(N,M)|
 \le\min\left\{L,\ L\left[2\zeta+\sqrt{\zeta^2+2c_N\zeta}\right]+\alpha\right\}.
\end{equation}
For $N=1$, $L(\delta_0+\eta)+\alpha$ is sufficient. The transfer preserves the same mean budget and total state set. Its square-root order in preparation error and in $(N-1)\eta$ is necessary in the full class, already for $M=2$, as proved by globally optimal singular-acquisition witnesses. No uniform absorption of the original competitors, normalized-filter contraction, or report recoverability is needed in this part. The bound is not asserted best possible in its numerical constants or sharp on every subclass.
\end{theorem}

Parts (i)--(iii), whose complete proofs are included, are the mean-acquisition implementation and geometric mechanisms developed in~\cite{GTFMean}. Part (iv), proved in full below, is the exact-budget constructive completion of~\cite{GTFExactBudget}. Part (v) is the budget-normalized transfer proved in Section~\ref{sec:sharp-transfer}; its all-program sharpness and the stronger memory-dependent singular frontier are proved in Section~\ref{sec:singular-memory}. Part (i) is the causal implementation theorem. Part (ii) uses a stronger report interface to turn implementation into a sharp acquired-geometry statement. Definition~\ref{def:recoverable} by itself does not require a finite affine closure; when both sets of hypotheses hold, the profile bounds are bounds on the exact value \eqref{eq:Phi}. A general normalized hidden-state filter need not satisfy recoverability. Part (iii) is not an assignment of dimension to arbitrary reachable states: its measure is an explicitly obtained reference law and its lower bound covers the actual stopped law of every competing policy.

\begin{remark}[What an occupation certificate certifies]\label{rem:ghost}
A candidate $z^\theta$ in \eqref{eq:budget-candidate} can contain invariant mass. The theorem certifies its terminal ensemble $d^\theta$, with an actual occupation no larger than $z^\theta$. It does not identify every candidate with the acquired occupation. This distinction is precisely what permits compactness without uniform absorption. Equality constraints on expected acquisition, rewards for spending resources, and negative running costs cannot be substituted in this argument.
\end{remark}

\begin{remark}[Fixed-time and stopping architectures]
The $M-1$ in \eqref{eq:exact-profile} pays for one live state. It is not the $M-n$ deduction for an internally counted, deterministic deadline. The implementing program uses no clock at all. It may run for arbitrarily many acquisitions, each charged in $\E\tau$. A comparison with a pathwise cap, a latency quantile, or a prescribed audit time requires a different resource constraint.
\end{remark}
